\section{The lower bound at a small absolute prefactor}
\label{sec:assembly}
We first prove Theorem~\ref{thm:main} for $Ns\le c_0$, where $c_0>0$ is an absolute constant chosen below. Section~\ref{sec:all-subcritical-prefactors} completes the remaining prefactor range. Fix once and for all
\[
 t=2^{-24},\qquad a=2^{-24},\qquad \Lambda=2^{16}.
\]
Use Theorem~\ref{thm:native-near-baseline} with output bound $A_0=1$ and a sufficiently small fixed $\delta>0$. Its input coefficient $\eta>0$ may be reduced without invalidating the construction. Set $D=\lfloor\eta H\rfloor$, $d=D+2$, and $n=M=4\cdot2^D$. Increasing the depth threshold ensures $D\ge1$.

Choose $\delta$ so small and then $H$ so large that the deterministic energy bound $K=(1+\delta)H$ satisfies
\begin{equation}
 q:=(1-a)(H+1)t^{2/(H+1)}\le K,
 \qquad 1-q/K\le2^{-18}.
 \label{eq:ratio-choice}
\end{equation}
This is possible because $q/H\to1-a$ and $K/H=1+\delta$. If needed enlarge $K$ slightly to ensure the first inequality; it already holds for all sufficiently large $H$ with the stated choices.

The empirical atom means $b(x)$ lie in $[-1,1]$. Assign arbitrary values in that interval at unobserved atoms. The native-small-noise memorizer, Theorem~\ref{thm:native-near-baseline}, supplies a Gaussian comparison law with mean energy at most $K$ and pointwise RMS error at most $2^{-10}$ against these means. Its expected empirical loss is at most $r_0+2^{-20}$: the within-atom residual has zero average at each atom, so its cross term with a function prediction vanishes exactly. Thus~\eqref{eq:memo-comparison} holds with $\zeta=2^{-20}$.

For $N\ge d$ and $H\ge2$, the operator moment constant obeys
\[
 \mathcal B_a(H,N,d)\le C_a HN,
 \qquad C_a=2\log(1/a)+8/a.
\]
Fix the absolute prefactor
\[
 c_0=\min\{2^{-24},2^{-22}/(\Lambda C_a)\}.
\]
At every $0<s\le c_0/N$ and every $\lambda\le\Lambda/H$,
\begin{equation}
 s+\lambda s\mathcal B_a(H,N,d)
 \le c_0/N+\Lambda C_a c_0\le2^{-20}.
 \label{eq:small-noise-choice}
\end{equation}
The native comparison exists uniformly throughout this scale range at every $N\ge e^{CH}$, for a sufficiently large absolute $C$. All choices are deterministic and independent of the sample.

For $\lambda\le\Lambda/H$, combine~\eqref{eq:finite-objective} with Lemma~\ref{lem:segment}, using the effective energy $(1-a)Q_{\mathrm{op}}$ and floor $q$. On the event~\eqref{eq:sample-event}, both the raw posterior mean and the mean of clipped posterior draws are within
\[
 u\le2^{-24}+\sqrt{2^{-18}+2^{-23}+2^{-20}+2^{-20}}<1/256
\]
of some vector $pb$, $p\in[0,1]$. Lemma~\ref{lem:witness} gives
\begin{equation}
 \langle b-2g(S),m(S)-g(S)\rangle_n
 \ge1/64-1/512-1/256>1/128.
 \label{eq:common-positive-witness}
\end{equation}
Clipping a raw mean also preserves its distance from $pb$, because $|b|\le1$.

For $\lambda\ge\Lambda/H$, use Lemma~\ref{lem:posterior-op-tail}. Since $\widehat R(0)\le1$ and $s\mathcal B_a\le C_aHc_0$, its moment bound has
\[
 \frac{q_*+2s(L/2+1)/(1-a)}L
 \le\frac{(1+c_0)/\Lambda+C_a c_0+3c_0}{1-a}<1/8.
\]
Thus every raw or clipped mean convention has pointwise absolute value at most $8^{-L/2}$, which is at most $2^{-15}$ after increasing the depth threshold. Directly,
\[
 \langle b-2g(S),m(S)-g(S)\rangle_n
 \ge\|g(S)\|_n^2-|\Delta|-\|b-2g(S)\|_n\|m(S)\|_n
 >1/128.
\]
This proves the same witness for every positive temperature and every allowed width and scale, on the one sample event.

The witness is linear in the prediction. Any data-dependent mixture of these mean predictions therefore obeys~\eqref{eq:common-positive-witness}. More strongly, every component is within $1/256$ of the segment $\{pb:0\le p\le1\}$; this remains true after mixing, and clipping the resulting mixture preserves that distance. Thus clipping after mixing is covered as well. The population conversion~\eqref{eq:population-witness}--\eqref{eq:excess-constant} gives excess at least $2^{-16}$.

The failure probabilities in~\eqref{eq:occupancy} and for the noise alignment tend to zero as $M=4\cdot2^{\lfloor\eta H\rfloor}$ grows. Increase the fixed depth threshold so their sum is at most $1/10$. This proves Theorem~\ref{thm:main} in the range $Ns\le c_0$.

\section{The fixed-width strong-prior limit}
For fixed $H,D,N$ above the deterministic interpolation threshold, choose the exact deterministic interpolator from Theorem~\ref{thm:near-baseline-memorizer} and keep its mean parameters fixed as $s\downarrow0$. The hidden-state perturbation estimate in Corollary~\ref{cor:near-baseline-native-small-noise} makes the prediction RMS of its native-covariance Gaussian perturbation tend to zero. The additive constant in its energy bound can be absorbed into the ratio choice~\eqref{eq:ratio-choice} by increasing the depth threshold. Also $s\mathcal B_a(H,N,d)\to0$. Hence~\eqref{eq:small-noise-choice} and the comparison error requirements hold for every sufficiently small positive $s$, uniformly over all temperatures $\lambda>0$.
The preceding proof therefore applies at every fixed width $N\ge e^{CH}$ throughout a sufficiently strong-prior interval $0<s<s_0(H,N)$. It applies in particular to every limit of the resulting predictions on the finite input support, including limits with data-dependent varying temperatures. No uniqueness of a parameter-space minimizer is assumed.

For completeness, for each fixed positive $\lambda$ the usual Gaussian Laplace argument concentrates parameter distributions on minimizers of $\widehat R_S+\lambda Q_N$. This objective is continuous and coercive. The objective-mixture argument already applies directly to any distribution on those minimizers. The theorem consequently does not depend on an arbitrary choice between tied minimizers.

